\section{Bloques completamente aleatorizados, cuadrados latinos y grecolatinos}

En innumerables escenarios prácticos, las unidades experimentales disponibles no son un conjunto estrictamente homogéneo. Por ejemplo, al evaluar motores de bases de datos en servidores físicos de diversas generaciones, o al probar fármacos en pacientes de distintos rangos de edad, una gran porción de la variabilidad observada en la variable respuesta no se debe al factor bajo estudio (el tratamiento), sino a las diferencias intrínsecas de las unidades experimentales.

Si realizáramos un ANOVA de un factor ignorando esa heterogeneidad conocida, dicha variación se absorbería directamente dentro de la Suma de Cuadrados del Error ($\text{SCE}$), inflando artificialmente el Cuadrado Medio del Error ($\text{CME}$) y reduciendo drásticamente la capacidad estadística del estadístico $F$ para detectar los efectos reales de los tratamientos.

El \textbf{Diseño en Bloques Completos al Azar (DBCA)} resuelve este problema agrupando las unidades experimentales en $b$ \textbf{bloques homogéneos} y aplicando los $k$ tratamientos aleatoriamente dentro de cada bloque, de modo que cada tratamiento aparece exactamente una vez por bloque ($N = k \times b$).

\subsection{Modelo aditivo y partición bidimensional de la varianza}

El modelo estadístico paramétrico para un diseño en bloques completos al azar se escribe como:
\begin{align}
	Y_{ij} = \mu + \tau_i + \beta_j + \epsilon_{ij}, \quad \text{para } i = 1, \ldots, k \text{ y } j = 1, \ldots, b,
	\label{eq:modelo_dbca}
\end{align}
donde:
\begin{itemize}
	\item $\mu$ es el promedio global de la población.
	\item $\tau_i$ es el efecto principal del \textbf{tratamiento $i$}, sujeto a la restricción $\sum_{i=1}^k \tau_i = 0$.
	\item $\beta_j$ es el efecto principal del \textbf{bloque $j$}, sujeto a la restricción $\sum_{j=1}^b \beta_j = 0$.
	\item $\epsilon_{ij} \sim N(0, \sigma^2)$ son los errores experimentales aleatorios i.i.d., asumiendo además la \textbf{ausencia de interacción} (aditividad pura) entre tratamientos y bloques.
\end{itemize}

\begin{teorema}[Teorema de partición de la suma de cuadrados en DBCA]
	En un Diseño en Bloques Completos al Azar con $k$ tratamientos y $b$ bloques, la Suma de Cuadrados Total se descompone en tres fuentes ortogonales de variación:
	\begin{align}
		\text{SCT} = \text{SCTR} + \text{SCB} + \text{SCE},
	\end{align}
	donde:
	\begin{itemize}
		\item $\text{SCTR} = b \sum_{i=1}^k (\bar{Y}_{i.} - \bar{Y}_{..})^2$ es la Suma de Cuadrados de los Tratamientos ($\nu_{\text{TR}} = k - 1$).
		\item $\text{SCB} = k \sum_{j=1}^b (\bar{Y}_{.j} - \bar{Y}_{..})^2$ es la \textbf{Suma de Cuadrados de los Bloques} ($\nu_{\text{B}} = b - 1$), donde $\bar{Y}_{.j} = \frac{1}{k}\sum_{i=1}^k Y_{ij}$ es el promedio observado de todo el bloque $j$.
		\item $\text{SCE} = \sum_{i=1}^k \sum_{j=1}^b (Y_{ij} - \bar{Y}_{i.} - \bar{Y}_{.j} + \bar{Y}_{..})^2$ es la Suma de Cuadrados del Error aleatorio residual, calculada en la práctica por diferencia algebraica directa:
		\begin{align}
			\text{SCE} = \text{SCT} - \text{SCTR} - \text{SCB},
		\end{align}
		con sus respectivos grados de libertad $\nu_E = (k - 1)(b - 1)$.
	\end{itemize}
\end{teorema}

El DBCA permite evaluar dos hipótesis nulas independientes: la prueba primordial sobre la igualdad de tratamientos ($H_{0,\text{TR}}: \tau_i = 0$) y la prueba secundaria para verificar si el factor de bloqueo redujo efectivamente el ruido del experimento ($H_{0,\text{B}}: \beta_j = 0$). Los estadísticos y su síntesis computacional se presentan en la Tabla \ref{tab:anova_dbca}.

\begin{table*}[htbp]
	\centering
	\caption{Tabla de Análisis de Varianza para un Diseño en Bloques Completos al Azar (DBCA).}
	\label{tab:anova_dbca}
	\begin{tabular}{lcccc}
		\hline
		\textbf{Fuente de Variación} & \textbf{SC} & \textbf{g.l. ($\nu$)} & \textbf{Cuadrado Medio (CM)} & \textbf{Estadístico $F$} \\ \hline
		Tratamientos & $\text{SCTR} = b \sum (\bar{Y}_{i.} - \bar{Y}_{..})^2$ & $k - 1$ & $\text{CMTR} = \frac{\text{SCTR}}{k - 1}$ & $F_{\text{TR}} = \frac{\text{CMTR}}{\text{CME}}$ \\ [1.5ex]
		Bloques & $\text{SCB} = k \sum (\bar{Y}_{.j} - \bar{Y}_{..})^2$ & $b - 1$ & $\text{CMB} = \frac{\text{SCB}}{b - 1}$ & $F_{\text{B}} = \frac{\text{CMB}}{\text{CME}}$ \\ [1.5ex]
		Error residual & $\text{SCE} = \text{SCT} - \text{SCTR} - \text{SCB}$ & $(k-1)(b-1)$ & $\text{CME} = \frac{\text{SCE}}{(k-1)(b-1)}$ & \\ [1.5ex] \hline
		\textbf{Total} & $\text{SCT} = \sum \sum (Y_{ij} - \bar{Y}_{..})^2$ & $kb - 1$ & & \\ \hline
	\end{tabular}
\end{table*}

\subsection{Cuadrados latinos}

El DBCA controla la variabilidad proveniente de \textbf{una sola} fuente de bloqueo conocida. Cuando existen \textbf{dos} fuentes de variabilidad extrañas conocidas y cruzadas entre sí (por ejemplo, el turno de trabajo y la máquina utilizada en una línea de producción), el \textbf{diseño de cuadrado latino} permite controlar ambas simultáneamente sin que el número de corridas experimentales crezca multiplicativamente.

\begin{definicion}[Cuadrado latino]
	Un \emph{cuadrado latino} de orden $n$ es un arreglo de $n \times n$ celdas en el que se distribuyen $n$ tratamientos (denotados tradicionalmente con letras latinas $A, B, C, \ldots$) de manera que cada tratamiento aparece \textbf{exactamente una vez en cada fila y exactamente una vez en cada columna}. Las filas y columnas representan los dos factores de bloqueo; el número de corridas experimentales es $n^2$ en lugar de las $n^3$ que requeriría un diseño factorial completo con tres factores de $n$ niveles cada uno.
\end{definicion}

\begin{center}
	\begin{tabular}{c|cccc}
		 & Columna 1 & Columna 2 & Columna 3 & Columna 4 \\ \hline
		Fila 1 & $A$ & $B$ & $C$ & $D$ \\
		Fila 2 & $B$ & $C$ & $D$ & $A$ \\
		Fila 3 & $C$ & $D$ & $A$ & $B$ \\
		Fila 4 & $D$ & $A$ & $B$ & $C$ \\
	\end{tabular}
\end{center}

El modelo estadístico paramétrico para un cuadrado latino de orden $n$ es
\begin{align}
	Y_{ijk} = \mu + \rho_i + \gamma_j + \tau_k + \epsilon_{ijk},
	\label{eq:modelo_cuadrado_latino}
\end{align}
donde $\rho_i$ es el efecto de la fila $i$, $\gamma_j$ el efecto de la columna $j$, y $\tau_k$ el efecto del tratamiento $k$ (con $i,j,k = 1,\ldots,n$, sujetos a las restricciones usuales $\sum \rho_i = \sum \gamma_j = \sum \tau_k = 0$); nótese que, a diferencia del DBCA, solo se observa \emph{una} de las $n^2$ combinaciones fila-columna-tratamiento posibles por cada par $(i,j)$, ya que el propio arreglo del cuadrado latino determina qué tratamiento $k$ corresponde a cada celda.

\begin{teorema}[Partición de la suma de cuadrados en un cuadrado latino]
	La Suma de Cuadrados Total se descompone en cuatro fuentes ortogonales de variación:
	\begin{align}
		\text{SCT} = \text{SC}_{\text{filas}} + \text{SC}_{\text{columnas}} + \text{SCTR} + \text{SCE},
	\end{align}
	con $n^2 - 1$ grados de libertad totales, repartidos como $\nu_{\text{filas}} = \nu_{\text{columnas}} = \nu_{\text{TR}} = n-1$ y $\nu_E = (n-1)(n-2)$ para el error residual.
\end{teorema}

\begin{center}
	\begin{tabular}{lcccc}
		\hline
		\textbf{Fuente de Variación} & \textbf{SC} & \textbf{g.l. ($\nu$)} & \textbf{Cuadrado Medio (CM)} & \textbf{Estadístico $F$} \\ \hline
		Filas & $\text{SC}_{\text{filas}}$ & $n-1$ & $\text{SC}_{\text{filas}}/(n-1)$ & --- \\ [1ex]
		Columnas & $\text{SC}_{\text{columnas}}$ & $n-1$ & $\text{SC}_{\text{columnas}}/(n-1)$ & --- \\ [1ex]
		Tratamientos & $\text{SCTR}$ & $n-1$ & $\text{CMTR} = \frac{\text{SCTR}}{n-1}$ & $F_{\text{TR}} = \frac{\text{CMTR}}{\text{CME}}$ \\ [1ex]
		Error residual & $\text{SCE}$ & $(n-1)(n-2)$ & $\text{CME} = \frac{\text{SCE}}{(n-1)(n-2)}$ & \\ [1ex] \hline
		\textbf{Total} & $\text{SCT}$ & $n^2 - 1$ & & \\ \hline
	\end{tabular}
\end{center}

\begin{observacion}
	Al igual que en el DBCA, la hipótesis de interés primordial es $H_{0,\text{TR}}: \tau_k = 0$ para todo $k$ (igualdad de tratamientos), contrastada mediante $F_{\text{TR}}$. Las filas y columnas son bloques de control, no variables de interés experimental directo, por lo que usualmente no se reportan sus estadísticos $F$.
\end{observacion}

\subsection{Cuadrados grecolatinos}

Cuando existe una \textbf{tercera} fuente de variabilidad extraña conocida (además de las dos ya controladas por filas y columnas), el \textbf{cuadrado grecolatino} extiende la idea del cuadrado latino superponiendo un segundo arreglo de letras (tradicionalmente griegas: $\alpha, \beta, \gamma, \ldots$) sobre el mismo cuadrado $n \times n$.

\begin{definicion}[Cuadrado grecolatino]
	Un \emph{cuadrado grecolatino} de orden $n$ consiste en dos cuadrados latinos \textbf{ortogonales} superpuestos: uno con letras latinas ($A, B, C, \ldots$, el factor de tratamiento principal) y otro con letras griegas ($\alpha, \beta, \gamma, \ldots$, un cuarto factor controlado), de modo que cada combinación (letra latina, letra griega) aparece exactamente una vez en todo el arreglo. Se dice que los dos cuadrados son \emph{ortogonales} porque cada par (letra latina, letra griega) coincide en exactamente una celda.
\end{definicion}

\begin{observacion}[Condición de existencia]
	No existen cuadrados grecolatinos ortogonales de orden $n=6$ (resultado célebre conjeturado por Euler en 1782 y probado en 1900); para todo otro $n \geq 3$ sí existen. En la práctica, los cuadrados grecolatinos se usan casi exclusivamente para $n \geq 3$, $n \neq 6$.
\end{observacion}

El modelo estadístico paramétrico añade un cuarto término al modelo del cuadrado latino:
\begin{align}
	Y_{ijkl} = \mu + \rho_i + \gamma_j + \tau_k + \delta_l + \epsilon_{ijkl},
	\label{eq:modelo_cuadrado_grecolatino}
\end{align}
donde $\delta_l$ es el efecto del factor griego $l$ (con $i,j,k,l = 1,\ldots,n$). La Suma de Cuadrados Total se descompone ahora en cinco fuentes ortogonales:
\begin{align}
	\text{SCT} = \text{SC}_{\text{filas}} + \text{SC}_{\text{columnas}} + \text{SCTR} + \text{SC}_{\text{griego}} + \text{SCE},
\end{align}
con $\nu_{\text{filas}} = \nu_{\text{columnas}} = \nu_{\text{TR}} = \nu_{\text{griego}} = n-1$ y $\nu_E = (n-1)(n-3)$ grados de libertad para el error residual, sobre un total de $n^2 - 1$.

\begin{observacion}
	Cada fuente de bloqueo adicional (fila, columna, factor griego) que efectivamente explica variabilidad real reduce el Cuadrado Medio del Error y, por lo tanto, aumenta la potencia de la prueba $F$ sobre los tratamientos —el mismo principio de bloqueo del DBCA, aplicado ahora a tres fuentes de ruido simultáneas en lugar de una sola.
\end{observacion}
